% ============================ Enrico Ribiani 16-03-2021 ====================================================================
% Base per i documenti  
\documentclass{article}
% ------------ pacchetti necessari ----------------
\usepackage[a4paper, total={6in, 8in}]{geometry} % formattazione decente della pagina
\usepackage{graphicx}                            % need for figure
\usepackage{amsmath}
\usepackage{amsfonts}                            % if you want the fonts
\usepackage{amssymb}                             % if you want extra symbols
\usepackage{graphicx}                            % need for figures
\usepackage{mathptmx}
\usepackage{float}                               % serve per mettere tabelle e immagini dove si vuole 
\usepackage[utf8]{inputenc}
\usepackage{textcomp}
\usepackage[hang,flushmargin,bottom]{footmisc}   % footnote format
\usepackage{fancyhdr, lastpage}
\usepackage{titlesec}
\titleformat{\section}{\normalsize\bfseries}{\thesection.}{1em}{}	% required for heading numbering style

\titleformat*{\section}{\Large\bfseries}
\titleformat*{\subsection}{\large\bfseries}
%\titleformat*{\subsubsection}{\large\bfseries}
%===================inizio pagina del titolo=================
\begin{document}
    \begin{titlepage}
    \begin{center}
% ------------------ inizio immagine logo ----------
\begin{figure}
    \centering
    \includegraphics{/home/rib/varie/logo.png}
    \label{fig:logo}
\end{figure}
% ------------------ fine immagine logo ----------
% ------------------ fine immagine logo ----------
-------------------------------------------------------------------------------------\\
\vspace{2\baselineskip}
\large Enrico Ribiani\\
\large 3AUB\\
\vfill

\Huge{\textbf{Riassunto Modulo 2, Unità 2-3-4}}\\
\vfill

%\LARGE{numero esperienza}\\
%\vfill
\LARGE{\textbf{19-04-2021}}\\
%\vfill
\end{center}
%=============== fine pagina titolo ===============
\end{titlepage}
\tableofcontents
\vspace{3\baselineskip}

\section{Unità 2: Resistori}

\subsection{Parametri e caratteristiche}
I resistori sono uno dei componenti passivi più frequenti all'interno di un circuito\\
Essi comportano l'introduzione di una \textbf{resistenza} elettrica nel circuito, misurata in ohm [$\Omega$].
I principali parametri ossia quelli più importanti, che vanno a influire maggiormente sulla scelta di un resistore sono:\\
\begin{enumerate}
    \item valore nominale
    \item tolleranza
    \item potenza nominale
\end{enumerate}
Il \textbf{valore nominale} rappresenta il valore più probabile a cui fare riferimento durante il progetto del circuito.\\
La \textbf{tolleranza} rappresenta il campo di valori entro cui il valore reale della resistenza è compreso, spesso si esprime in maniera relativa al valore nominale.\\
La \textbf{potenza nominale} indica la massima potenza dissipabile a temperatura ambiente, normalmente \textit{25 °C}, senza che la sua resistenza venga alterata.
Le resistenze sono classificate in serie dalle \textbf{serie} che raggruppano dei valori commerciali, questa parte non verrà riassunta.\\
\vspace{1\baselineskip} 
Il valore nominale dei resistori e la tolleranza si può leggere grazie a un sistema di bande colorate poste sulle resistenze associate a dei valori tramite il codice colori.\\
Ci sono tre tipi di codici colori, quello a 4 bande, quello a 5 e quello a 6; tutti e tre seguono lo schema: \emph{Cifre significative, moltiplicatore, tolleranza}, infatti le prime 3 bande (prendendo per esempio i codici a 5 bande) indicano le cifre significative che indicano le prime cifre del valore della resistenza, la $4^a$ cifra indica l'ordine di grandezza o moltiplicatore $10^n$ mentre l'ultima ne indica la tolleranza, nei codici a 6 bande viene indicata dalla $6^a$ banda il coefficiente di temperatura che indica come varia la resistenza in base alla temperatura.
\vspace{\baselineskip}[H]
Una caratteristica intrinseca molto importante è la  \textbf{potenza dissipabile}, essa indica la potenza massima che il resistore è in grado di dissipare senza che si rovini ossia senza raggiungere valori troppo alti dal momento che l'energia dissipata viene trasformata in calore.\\
Per calcolarla bisogna tenere conto che se il resistore è più caldo della temperatura ambiente parte del calore viene assorbito \textit{potenza dissipata}.
\begin{center}
    La cui formula è: $P_d=\frac{T_i-T_a}{R_th}$.\\
$T_i$ = temperatura interna del componente, $T_a$ = temperatura ambiente, $R_th$=resistenza termica.
\end{center}
Ovviamente la potenza dissipabile è inversamente proporzionale alla temperatura ambiente.\\
Questa variazione è prevedibile e calcolabile grazie al coefficiente di temperatura (lettera $\alpha$) e dalla formula:
\begin{center}
    $R=R_0\cdot[1+\alpha\cdot(T-T_0)]$
\end{center}
Allo stesso modo oltre alla resistenza de viene influenzata la resistività.\\
%___________________________----------------=====================================----------------___________________________do di
\subsection{Tecnologie costruttive}
La parti che formano un resistore sono:
\begin{itemize}
    \item elemento resistivo
    \item supporto dell'elemento resistivo
    \item rivestimento di protezione
    \item reofori
\end{itemize}
I vari tipi di resistori si distinguono per come l'\textbf{elemento resistivo}, ossia la parte che viene attraversata dalla corrente, viene utilizzato.\\
Ad esempio nei \textbf{resistori a impasto} l'elemento resistivo è composto da un agglomerato di materiali e \emph{polvere di carbone}\\
Nei \textbf{resistori a film} è costituito da uno strato di conduttore avvolto su uno strato di isolante, per aumentare la resistenza si riduce la sezione e si aumenta la resistenza dal momento che $R=\rho\cdot\frac{l}{S}$\\
Mentre nei\textbf{ resistori a filo} l'elemento resistivo è un filo metallico avvolto su materiale isolante con un "doppio avvolgimento" o \textbf{avvolgimento antinduttivo} per evitare fenomeni elettromagnetici.\\
Il materiale con cui viene fatto \textbf{supporto dell'elemento resistivo} deve avere una resistenza di isolamento elevata e proprietà meccaniche e termiche sufficienti a resistere alle temperature.\\
Rivestimento, resine o smalti.
%___________________________----------------=====================================----------------___________________________
\subsection{Tipi di resistori}
Ci sono due modi in cui i resistori si collegano ai circuiti: tramite reofori o THT \textit{Through Hole Technology} oppuresaldati al circuito o SMT \textit{Sourface Mounted Technology}.
\vspace{1\baselineskip}
\\Ci sono anche dei componenti che variando la resistenza riescono a modulare un valore come:
\begin{itemize}
    \item \textbf{trimmer}\\
    I trimmer sono resistenze variabili tramite una vite che scorre sull'elemento resistivo, non sono concepiti per essere variati molto.
    \item \textbf{potenziometri}\\
    i potenziometri vengono utilizzati per variare la tensione tramite partitore di tensione.
    \item \textbf{reostati}\\
    i reostati vengono usati per regolare la corrente. $I=\frac{V_g}{R_v+R_u}$ dove $Rv$ è la resistenza del trimmer.
\end{itemize}

%___________________________----------------=====================================----------------___________________________
\section{Unità 3: Condensatori}
\subsection{Generalità}
Un condensatore è formato da due armature separato da un dielettrico e ha la finzione di immagazzinare carica elettrica per poi cederla, quando è carico alle sue due stremità ci sarà una \textit{ddp} che dipende dalla quantità di \textbf{carica} e dalla \textbf{capacità}: $V_c=\frac{Q}{C}$\\
La carica di un condensatore p definibile come "L'attitudine di immagazzinare carica elettrica".\\
Se il condensatore è alimentato da un generatore di tensione si caricherà istantaneamente mentre con una resistenza in serie si caricherà esponenzialmente in $4,6\cdot\tau$ e $\tau=Req\cdot C$.\\
La formula generale della tensione o della carica di un condensatore è:
\begin{center}
    $V_{(t)}=V_{0}(1-e)^{-t/\tau}$\\
    $Q_{(t)}=Q_{0}(1-e)^{-t/\tau}$
\end{center}

\subsection{Parametri}
I principali parametri che caratterizzano un resistore sono:
\begin{enumerate}
    \item \textbf{capacità nominale}\\
    È il valore capacitivo determinato dal produttore, si misura in farad [F] e varia in base al dielettrico e alla geometria.\\
    Per avere una grande capacità è necessario avere un dielettrico molto sottile e con molta superficie oltre a una elevata costante dielettrica poiché $C=\epsilon\cdot\frac{S}{d}$
    \item \textbf{tolleranza}\\
    Analoga alle resistenze tranne per il fatto che può essere anche asimmetrica ad esempio $ -20\%\div+50\%$\\
    \item \textbf{coefficiente di temperatura}\\
    Indica il variare della capacità in base alla temperatura calcolabile in modo analogo alle resistenze:
    \begin{center}
        $C=C_0\cdot[1+T_c\cdot(T-T_0)]$\\
        $T_c$= Temperature Coefficient
    \end{center}

    \item \textbf{tensione nominale }\\
    È la massima tensione sopportabile dal dielettrico prima di bucarsi, quindi possiamo dedurre che questa grandezza sia inversamente proporzionale alla \emph{capacità}
    \item \textbf{resistenza di isolamento}\\
    La resistenza del dielettrico, si calcola analogamente a quella dei resistori $R=\rho\cdot\frac{d}{S}$
    \item \textbf{angolo di perdita}\\
    Rappresenta la perdita di energia dato dalle resistenze che dissipano energia per l'effetto Joule, la \emph{tangente} di questo angolo rappresenta il rapporto tra potenza persa e potenza reattiva cioè il \emph{fattore di dissipazione}.
\end{enumerate}
Alcuni di questi parametri sono riportati sui condensatori più grossi o nei \textit{datasheet}. Nel primo caso vengono riportate tensione nominale, capacità e tolleranza tramite lettere o codici alfanumerici.
\subsection{Tecnologie costruttive}
Ci sono 4 tipi di condensatori:
\begin{itemize}
    \item \textbf{condensatori ceramici}\\
    Chiamati così perché appunto usano la  ceramica come dielettrico, sono a forma di disco e di tubetto e vanno da pochi pF a qualche nF.\\
    Sono divisi in 3 classi, quelli della \textit{I classe} sono \textit{molto stabili} a livello di temperatura. 
    Quelli di \textit{II classe} hanno un comportamento \textit{poco lineare} ma hanno una $\epsilon$ maggiore di quelli di classe I e posso essere stabili o instabili.
    Mentre i condensatori di \textit{III classe} sono caratterizzati dallo spessore estremamente \textit{sottile} del \textit{dielettrico} per ridurre al massimo le dimensioni, ma di conseguenza la tensione nominale avrà un valore molto basso.
    
    \item \textbf{condensatori a film plastici}\\
   Questi condensatori usano come dielettrico un film di \textit{plastica metallizzata} come polistirolo, teflon ecc...\\
   Vengono utilizzati a bassa frequenza, hanno un valore che va da 1 nF a  1000nF e sono molto stabili in perdite e temperatura.
    
    \item \textbf{condensatori elettrolitici}\\
    Utilizzano come dielettrico un ossido di metallo, utilizzano come armature il metallo (alluminio o tantalio);i condensatori in alluminio usano $AI_2O_3$ come ossido.\\
    Il dielettrico è particolarmente sottile perché viene prodotto tramite elettrolisi utilizzando anche un foglio di rame, l'ossido deve venire riprodotto quando è in funzione collegando una delle due armature con l'elettrolita al + altrimenti l'ossido sparisce provocando una corrente intensa e la seguente esplosione data dall'aumento di temperatura; si dice che i condensatori sono polarizzati.\\
    I condensatori al tantalio si differenziano da quelli in alluminio solo per la porosità che permette una costruzione di foma a goccia e a parallelepipedo oltre al cilindro. Coprono valori da 1$\mu$F a 100mF, hanno valori nominali basse e hanno una piccola resistenza di isolamento.
    
    \item \textbf{condensatori variabili}\\
    Usano come dielettrico l'aria e sono formati da lamine di alluminio disposte a libro collegate in modo alternato,di due gruppi di armature uno è fisso e l'altro mobile, consentendo di variare la capacità cambiando la superficie in comune tra due armature.
    \begin{figure}[H]
        \centering
        \includegraphics[scale=0.2]{/home/rib/varie/capscaps.jpg}
        \caption{Condensatore variabile}
        \label{varcap}
    \end{figure}
\end{itemize}
\section{Unità 4: Induttori}
\subsection{Generalità}
\subsection{Parametri}
\subsection{Tecnologie costruttive}

\section{Unità 4: Relè}
\subsection{Generalità}
\subsection{Parametri}
\subsection{Tecnologie costruttive}


\end{document}